% Shared, project-local BZ kinetics infrastructure.
% Its source content is the T2 mechanism already transcribed in the T2 chapters.
\section{Shared Belousov--Zhabotinsky kinetics}

\subsection{Formalization contract and informal proof}

This project-local module factors only the numerical species, state, kinetic
constant, and mass-action-rate vocabulary common to T2-A2, T2-A3, and T2-A5.
It contains no target-specific stationary-state, period, or numerical-answer
theorem.  Its three generic algebraic lemmas only cancel shared positive
factors or specialize a state update.  All concentrations and rates are scalar
numerical readouts in the units printed by the source.

\begin{definition}

[Molar concentration]
\label{def:chem:bz:molar_concentration}
\lean{IChO2026Chem.Kinetics.BelousovZhabotinsky.MolarConcentration}
A molar concentration is represented by its real numerical value in the
source's molar scale.
\end{definition}
\begin{proof}

The common T2 calculations use real arithmetic on numerical concentration
readouts.
\end{proof}

\begin{definition}

[Second-order rate constant]
\label{def:chem:bz:second_order_rate_constant}
\lean{IChO2026Chem.Kinetics.BelousovZhabotinsky.SecondOrderRateConstant}
A second-order rate constant is represented by its real numerical value in
the printed $\mathrm{M^{-1}s^{-1}}$ scale.
\end{definition}
\begin{proof}

This type records the units of the constants multiplying two concentration
factors in the shared mass-action model.
\end{proof}

\begin{definition}

[Third-order rate constant]
\label{def:chem:bz:third_order_rate_constant}
\lean{IChO2026Chem.Kinetics.BelousovZhabotinsky.ThirdOrderRateConstant}
A third-order rate constant is represented by its real numerical value in
the printed $\mathrm{M^{-2}s^{-1}}$ scale.
\end{definition}
\begin{proof}

This type records the units of the constants multiplying three concentration
factors in the shared mass-action model.
\end{proof}

\begin{definition}

[Fourth-order rate constant]
\label{def:chem:bz:fourth_order_rate_constant}
\lean{IChO2026Chem.Kinetics.BelousovZhabotinsky.FourthOrderRateConstant}
A fourth-order rate constant is represented by its real numerical value in
the printed $\mathrm{M^{-3}s^{-1}}$ scale.
\end{definition}
\begin{proof}

This type records the units of the step-(5) constant multiplying four
concentration factors.
\end{proof}

\begin{definition}

[BZ species]
\label{def:chem:bz:species}
\lean{IChO2026Chem.Kinetics.BelousovZhabotinsky.Species}
The finite species type includes HBrO$_2$, BrO$_3^-$, Br$^-$, H$^+$,
BrO$_2^\bullet$, HBrO, Ce(III), Ce(IV), BMA, and malonic acid.
\end{definition}
\begin{proof}

These are the species named in the elementary-step mechanism used by the
three T2 targets.
\end{proof}

\begin{definition}

[BZ concentration state]
\label{def:chem:bz:state}
\lean{IChO2026Chem.Kinetics.BelousovZhabotinsky.State}
\uses{def:chem:bz:molar_concentration, def:chem:bz:species}
A state assigns a molar concentration to each shared BZ species.
\end{definition}
\begin{proof}

The mass-action expressions require only lookup of the current concentration
of each reactant.
\end{proof}

\begin{definition}

[BZ kinetic parameters]
\label{def:chem:bz:kinetic_parameters}
\lean{IChO2026Chem.Kinetics.BelousovZhabotinsky.KineticParameters}
\uses{def:chem:bz:second_order_rate_constant, def:chem:bz:third_order_rate_constant, def:chem:bz:fourth_order_rate_constant}
The parameter record supplies $k_1,\ldots,k_7$, typed by the order of their
corresponding elementary reactions.
\end{definition}
\begin{proof}

The fields are the source-given kinetic constants, kept separate from a state
and from all requested numerical conclusions.
\end{proof}

\begin{definition}

[Elementary rate 1]
\label{def:chem:bz:rate1}
\lean{IChO2026Chem.Kinetics.BelousovZhabotinsky.rate1}
\uses{def:chem:bz:kinetic_parameters, def:chem:bz:state}
The step-(1) rate is $k_1[\mathrm{HBrO_2}][\mathrm{BrO_3^-}][\mathrm{H^+}]$.
\end{definition}
\begin{proof}

It is the mass-action product for the three reactants in elementary step
(1).
\end{proof}

\begin{definition}

[Elementary rate 2]
\label{def:chem:bz:rate2}
\lean{IChO2026Chem.Kinetics.BelousovZhabotinsky.rate2}
\uses{def:chem:bz:kinetic_parameters, def:chem:bz:state}
The step-(2) rate is
$k_2[\mathrm{BrO_2^\bullet}][\mathrm{Ce^{III}}][\mathrm{H^+}]$.
\end{definition}
\begin{proof}

It is the mass-action product for the three reactants in elementary step
(2).
\end{proof}

\begin{definition}

[Elementary rate 3]
\label{def:chem:bz:rate3}
\lean{IChO2026Chem.Kinetics.BelousovZhabotinsky.rate3}
\uses{def:chem:bz:kinetic_parameters, def:chem:bz:state}
The step-(3) rate is $k_3[\mathrm{HBrO_2}]^2$.
\end{definition}
\begin{proof}

Two HBrO$_2$ reactants give the squared concentration factor in the
mass-action expression.
\end{proof}

\begin{definition}

[Elementary rate 4]
\label{def:chem:bz:rate4}
\lean{IChO2026Chem.Kinetics.BelousovZhabotinsky.rate4}
\uses{def:chem:bz:kinetic_parameters, def:chem:bz:state}
The step-(4) rate is $k_4[\mathrm{HBrO_2}][\mathrm{Br^-}][\mathrm{H^+}]$.
\end{definition}
\begin{proof}

It is the mass-action product for the three reactants in elementary step
(4).
\end{proof}

\begin{definition}

[Elementary rate 5]
\label{def:chem:bz:rate5}
\lean{IChO2026Chem.Kinetics.BelousovZhabotinsky.rate5}
\uses{def:chem:bz:kinetic_parameters, def:chem:bz:state}
The step-(5) rate is
$k_5[\mathrm{BrO_3^-}][\mathrm{Br^-}][\mathrm{H^+}]^2$.
\end{definition}
\begin{proof}

The two proton reactants supply the squared proton-concentration factor in
the fourth-order mass-action product.
\end{proof}

\begin{definition}

[Elementary rate 6]
\label{def:chem:bz:rate6}
\lean{IChO2026Chem.Kinetics.BelousovZhabotinsky.rate6}
\uses{def:chem:bz:kinetic_parameters, def:chem:bz:state}
The step-(6) rate is
$k_6[\mathrm{HBrO}][\mathrm{MA}]$.
\end{definition}
\begin{proof}

Step (6) has two reactants, HBrO and malonic acid, so its mass-action rate is
the second-order constant times their two concentration readouts.
\end{proof}

\begin{definition}

[Elementary rate 7]
\label{def:chem:bz:rate7}
\lean{IChO2026Chem.Kinetics.BelousovZhabotinsky.rate7}
\uses{def:chem:bz:kinetic_parameters, def:chem:bz:state}
The step-(7) rate is $k_7[\mathrm{Ce^{IV}}][\mathrm{BMA}]$.
\end{definition}
\begin{proof}

It is the mass-action product for the two reactants in the continuous
Process-C elementary step.
\end{proof}

\begin{lemma}

[Cancellation in the step-(1)/step-(4) switching balance]
\label{lem:chem:bz:rate1_eq_rate4_of_positive}
\lean{IChO2026Chem.Kinetics.BelousovZhabotinsky.rate1_eq_rate4_of_positive}
\uses{def:chem:bz:rate1, def:chem:bz:rate4}
If the step-(1) and step-(4) rates are equal and the common
$[\mathrm{HBrO_2}]$ and $[\mathrm{H^+}]$ factors are positive, then
$k_1[\mathrm{BrO_3^-}]=k_4[\mathrm{Br^-}]$.
\end{lemma}
\begin{proof}

Reassociate both mass-action products so that the two positive common factors
form one nonzero left factor, then cancel that factor.
\end{proof}

\begin{lemma}

[Step-(4) rate after updating bromide]
\label{lem:chem:bz:rate4_update_bromide}
\lean{IChO2026Chem.Kinetics.BelousovZhabotinsky.rate4_update_bromide}
\uses{def:chem:bz:rate4}
Updating only the bromide coordinate specializes the step-(4) rate to
$k_4[\mathrm{HBrO_2}][\mathrm{Br^-}][\mathrm{H^+}]$ with the supplied new
bromide value.
\end{lemma}
\begin{proof}

The updated bromide lookup reduces by decidable equality, while every other
species lookup remains the original state value.
\end{proof}

\begin{lemma}

[Step-(5) rate after updating bromide]
\label{lem:chem:bz:rate5_update_bromide}
\lean{IChO2026Chem.Kinetics.BelousovZhabotinsky.rate5_update_bromide}
\uses{def:chem:bz:rate5}
Updating only the bromide coordinate specializes the step-(5) rate to
$k_5[\mathrm{BrO_3^-}][\mathrm{Br^-}][\mathrm{H^+}]^2$ with the supplied new
bromide value.
\end{lemma}
\begin{proof}

As for step (4), simplification of the coordinate update yields the stated
mass-action product.
\end{proof}
